\paragraph{Scattering vs black-hole queue equivalence (toy).} \AuditTag We compare (i) a delay-driven scattering queue (two-channel $S(\omega)$ proxy, with Wigner--Smith delay setting a tick-time service cost) to (ii) a black-hole-like internal queue emitting one stable-type record $w_t\in X_m$ per tick. The shared control parameter is the periodic arrival rate (one job per $p$ ticks). Trap-like saturation is audited by the fraction of ticks where the internal queue length exceeds a fixed threshold $q_\star$, and by the maximum observed queue length. External-record structure is summarized by a simple compressibility proxy (zlib compression ratio) on the emitted record stream.
\paragraph{Scope.} \AuditTag This is an interface-level comparison artifact (CS model). It does not assert a theorem-level equivalence between scattering and black-hole physics.
\paragraph{Determinism.} \AuditTag Parameters are fixed (ticks=6000, q\_star=50, m=6, grid=81). BH knobs are CAP-selected from a bounded family (absorb\_batch in [1, 2, 3, 4], leak\_amp in [0.0, 0.05, 0.091, 0.1, 0.2, 0.35]) to match the lossy-scattering audit targets across arrival periods (RMSE=0.268970, gap=0.000092); selected: absorb\_batch=2, leak\_amp=0.091. Loss is reported via an S1-style $\eta_{\min}$ gate and an absorption fraction proxy in the lossy scattering case.